\documentclass[12pt,a4paper]{article}
\usepackage[utf8]{inputenc}
\usepackage[T1]{fontenc}
\usepackage[spanish]{babel}
\usepackage{geometry}
\usepackage{hyperref}
\usepackage{listings}
\usepackage{booktabs}
\usepackage{graphicx}
\geometry{margin=2.5cm}

\title{Informe t\'ecnico\\Sistema de gesti\'on de inventario\\Mercado Municipal de Quevedo}
\author{Zamora Arias Carla Esthefania\\Ingenier\'ia de Software --- UTEQ\\Aplicaciones Web}
\date{\today}

\begin{document}
\maketitle

\begin{abstract}
Se presenta la implementaci\'on de una API REST con Spring Boot 3 (Java 21) para el control de productos del Mercado Municipal de Quevedo. El servicio expone recursos paginados con contrato JSON uniforme, validaci\'on de entrada, borrado l\'ogico, cach\'e cache-aside con Redis y seguridad basada en JWT.
\end{abstract}

\section{Introducci\'on}
El m\'odulo de inventario digitaliza el control de productos mediante una arquitectura por capas (controlador, servicio, repositorio), alineada a las Unidades III y IV de la asignatura \cite{walls2022,martin2017}.

\section{Arquitectura}
\begin{itemize}
  \item \textbf{Controlador:} endpoints REST delgados.
  \item \textbf{Servicio:} reglas de negocio, cach\'e e invalidaci\'on.
  \item \textbf{Repositorio:} acceso exclusivo a PostgreSQL con Spring Data JPA.
\end{itemize}

\section{Modelo de datos}
La entidad \texttt{Producto} incluye \texttt{id}, \texttt{nombre}, \texttt{categoria}, \texttt{stock}, \texttt{precio}, \texttt{activo} y \texttt{creadoEn} (\texttt{TIMESTAMPTZ}). La eliminaci\'on es l\'ogica (\texttt{activo=false}).

\section{Requisitos implementados}
\begin{enumerate}
  \item \texttt{GET /api/v1/productos} paginado con envoltura \texttt{\{success, data, message, meta\}}.
  \item \texttt{POST /api/v1/productos} con \texttt{@Valid} y respuesta 400 por campo.
  \item \texttt{DELETE /api/v1/productos/\{id\}} con soft delete y 404.
  \item Cache-aside Redis: \texttt{@Cacheable} / \texttt{@CacheEvict} \cite{kleppmann2017}.
  \item JWT: \texttt{ROLE\_USER} (GET), \texttt{ROLE\_ADMIN} (POST/DELETE), c\'odigos 401/403 \cite{rfc7519}.
\end{enumerate}

\section{Seguridad}
Se usa un filtro Bearer que valida el token. El login (\texttt{POST /api/v1/auth/login}) emite JWT firmado. El secreto se configura por variable de entorno \texttt{JWT\_SECRET} y no se versiona en claro.

\section{Evidencias de prueba}
Las pruebas se documentan en \texttt{docs/requests.http} y capturas en \texttt{docs/capturas/} (200, 400, 401, 403, 404, cache hit/miss).

\section{Conclusiones}
La API cumple el contrato del examen parcial: es reproducible, segura con JWT y acelera lecturas con Redis, consolidando persistencia, patrones y servicios REST.

\bibliographystyle{plain}
\bibliography{referencias}
\end{document}
